\documentclass[10pt,a4paper]{amsart}
\usepackage[margin=19mm]{geometry}
\usepackage{amsmath,amssymb,mathtools}
\usepackage[scr=rsfs,cal=euler]{mathalfa}
\usepackage{xfrac}
\usepackage[unicode,hidelinks,bookmarks=false]{hyperref}
\newcommand{\ie}{i.e.\ }
\newcommand{\cf}{cf.\ }
\newcommand{\resp}{respectively\ }
\newcommand{\Subsection}{\subsection}
\providecommand{\nobreakdash}{\nobreak\hbox{-}}
\setlength{\emergencystretch}{2em}
\multlinegap0pt
\usepackage{bm}
\usepackage{bbm}
\usepackage{dsfont}
\usepackage{xspace}
\usepackage{cancel}
\usepackage{mathtools}
\usepackage{amsthm}
\makeatletter
\@ifundefined{assumption}{\newtheorem{assumption}{Assumption}}{}
\@ifundefined{cor}{\newtheorem{cor}{Corollary}}{}
\@ifundefined{lemma}{\newtheorem{lemma}{Lemma}}{}
\@ifundefined{thm}{\newtheorem{thm}{Theorem}}{}
\makeatother
\title[Nash Equilibria of Noncooperative/Mixed Differential Games w]{Nash Equilibria of Noncooperative/Mixed Differential Games with Density Constraints in Infinite Dimensions}
\author{Zhun Gou, Nan-Jing Huang, Jian-Hao Kang, Jen-Chih Yao}
\date{Source: arXiv:2508.07623v2 (CC BY 4.0)}
\begin{document}
\maketitle

\noindent\emph{Source and licence.} Nash Equilibria of Noncooperative/Mixed Differential Games with Density Constraints in Infinite Dimensions. Version arXiv:2508.07623v2 (CC BY 4.0), distributed under the Creative Commons Attribution 4.0 International licence, as recorded in the arXiv licence metadata for this exact version and shown on the abstract page https://arxiv.org/abs/2508.07623v2. Full rights notice: licenses/arxiv-2508.07623-CC-BY-4.0.txt.\par\smallskip

\noindent\textit{Editorial scope.} Counted unit: the Thm whose source label is \texttt{47} in the article (source0.tex lines 704--706, source label \texttt{47}), delivered together with the article's own complete original proof of it (source0.tex lines 707--784, one \texttt{proof} environment of 78 lines). No mathematics is written, repaired or re-derived: statement and proof are the article's own text line for line. Required context retained uncounted: 15 (lines 316--327); 17 (lines 360--362); 16 (lines 365--368); 29 (lines 371--373); 35 (lines 576--578); 21 (lines 609--619); 33 (lines 647--649); 37 (lines 684--686). Every definition, display and label of those lines is retained unchanged. Omitted from this package and not redistributed: 1--315, 328--359, 363--364, 369--370, 374--575, 579--608, 620--646, 650--683, 687--703, 785--1021. Nothing inside a retained excerpt is omitted or altered: every retained body is delivered exactly as the article's lines are written, so no omission receipt of any kind accompanies this package. Citations: the retained excerpts carry no citation command at all, so no bibliography entry of the article is transcribed and no reference is redistributed. Reference adaptation: none was needed and none was made; every reference inside the delivered unit resolves inside it, so no unresolved reference or citation is produced. Printed slips kept literal: no printed word, symbol, hypothesis or number of the counted statement or of its proof was altered. Layout and class: the article's own class is \texttt{article} with packages including latexsym, amsfonts, amsthm, amsmath, amssymb, cases, authblk, multicol, chngcntr, graphicx, algorithm2e, hyperref, enumerate, epstopdf; the wrapper is the lane's pinned \texttt{amsart} editorial wrapper at 10pt on A4, which restates the theorem-like environments the retained text needs and adds the article's own macro definitions that the retained text needs (none), with extra wrapper packages (bm, bbm, dsfont, xspace, cancel, mathtools). The delivered numbering is the unit's own. Assets: no retained span contains a figure environment, a graphics-inclusion command, a PDF-page inclusion, a table or a TikZ picture; no image file is copied into this package, the package declares no asset dependency and no image file is redistributed with it. Licence: version arXiv:2508.07623v2 of this article is distributed under the Creative Commons Attribution 4.0 International licence, as recorded in the arXiv licence metadata for this exact version and shown on the abstract page https://arxiv.org/abs/2508.07623v2. A full rights notice is delivered at licenses/arxiv-2508.07623-CC-BY-4.0.txt. Characters: every retained excerpt is pure ASCII, so the delivered engine, XeLaTeX with its default Latin Modern text font, reports no missing character.
\begin{equation}\label{15}
\begin{cases}
\widehat{V}=P\widehat{u}+Q=\int_s^TP(t)\widehat{u}dt+Q,\quad \mathcal{P}_i(t)=Y_i(t)B_{ii}(t)+F_i(t),\\
\overline{Y}_i=\alpha_i+Y_{i}(s)\xi_{i}+\int_s^TY_{i}(t)f_{i}(t)dt,\quad Q=(\overline{Y}_1,\cdots,\overline{Y}_n)^{\top},\\
P(t)=\begin{bmatrix}
\mathcal{P}_1(t)& Y_1(t)B_{12}(t)&\dots & Y_1(t)B_{1n}(t)\\
  Y_2(t)B_{21}(t)& \mathcal{P}_2(t)&\dots & Y_2(t)B_{2n}(t)\\
\vdots &\vdots & \ddots&\vdots\\
Y_n(t)B_{n1}(t)&Y_n(t)B_{n2}(t)&\dots & \mathcal{P}_n(t)
\end{bmatrix}_{n \times n}.
\end{cases}
\end{equation}

\begin{equation}\label{17}
\widehat{u}=\mathcal{P}_{U^n}[\widehat{u}-\epsilon_0(P\widehat{u}+Q)]=\mathcal{P}_{U^n}[(\mathbb{I}_{H^n}-\epsilon_0P)\widehat{u}-\epsilon_0Q],
\end{equation}

\begin{assumption}\label{16}
Assume that $P\in \mathcal{L}(H^n)$ defined in \eqref{15} is positive definite, i.e., there exists a constant $\epsilon_P>0$ such that
$\langle PZ,Z\rangle_{H^n}\geq \epsilon_{P}\|Z\|^2_{H^n}$ for all $Z\in H^{n}$.
\end{assumption}

\begin{thm}\label{29}
Under Assumption \ref{16}, there exists a unique solution $\widehat{u}\in U^n$ to \eqref{17} when $\epsilon_0\in(0,2\epsilon_P\|P\|_{\mathcal{L}(H^n)}^{-2})$.
\end{thm}

\begin{equation}\label{35}
\langle P_k\widehat{v}_k+Q_k,Z-\widehat{v}_k\rangle_{H^n}\geq 0,\quad \forall Z\in U^n_{\varepsilon}.
\end{equation}

\begin{assumption}\label{21}
Let $B=[B_{ij}]_{i,j\in N}$ and $B_k=[B_{ijk}]_{i,j\in N}$. Assume that, for each $i\in N$ we have $\min_{k}\varepsilon_{1k}>-1$ and
\begin{equation*}
  \begin{cases}
  \lim_{k\rightarrow+\infty}\varepsilon_{1k}=\lim_{k\rightarrow+\infty}\varepsilon_{2k}^{(i)}=0,\quad\lim_{k\rightarrow+\infty}\|G_i-G_{ik}\|_{\mathcal{L}(H)}=0,\\
    \lim_{k\rightarrow+\infty}\|A_i-A_{ik}\|_{ C([s,T],\mathcal{L}(H))}=\lim_{k\rightarrow+\infty}\|B-B_{k}\|_{ C([s,T],\mathcal{L}(H^n))}=0, \\
  \lim_{k\rightarrow+\infty}\|\xi_i-\xi_{ik}\|_H=\lim_{k\rightarrow+\infty}\|f_i-f_{ik}\|_{H}=\lim_{k\rightarrow+\infty}\|\alpha_i-\alpha_{ik}\|_{H}=0, \\
  \lim_{k\rightarrow+\infty}\|E_i-E_{ik}\|_{ C([s,T],\mathcal{L}(H))}=\lim_{k\rightarrow+\infty}\|F_i-F_{ik}\|_{ C([s,T],\mathcal{L}(H))}=0.
  \end{cases}
\end{equation*}
\end{assumption}

\begin{lemma}\label{33}
Under Assumption \ref{21}, one has $$\lim_{k\rightarrow+\infty}\|P-P_k\|_{\mathcal{L}(H^n)}=\lim_{k\rightarrow+\infty}\|Q-Q_k\|_{H^n}=0.$$
\end{lemma}

\begin{cor}\label{37}
Under Assumptions \ref{16} and \ref{21},  there exists a positive integer  $\mathcal{N}(\epsilon_P)$ such that $P_k$ satisfies  Assumption \ref{16} for any $k>\mathcal{N}(\epsilon_P)$.
\end{cor}

\begin{thm}\label{47}
Under Assumptions \ref{16},  both the NNE  and the OWAP are stable.
\end{thm}

\begin{proof}
Let
\begin{equation}\label{50}
\varepsilon_{2k}=(\varepsilon^{(1)}_{2k},\cdots,\varepsilon^{(n)}_{2k})^{\top}, \quad Z_k=h(Z)=\frac{Z-\varepsilon_{2k}}{1+\varepsilon_{1k}},\quad \overline{v}_k=h(\widehat{v}_k)=\frac{\widehat{v}_k-\varepsilon_{2k}}{1+\varepsilon_{1k}},
\end{equation}
where the mapping $h:U^n_{\varepsilon}\rightarrow U^n$ is  bijective, then one can check that $Z_k,\overline{v}_k\in U^n$ for any $k$, and
$\lim_{k\rightarrow+\infty}\|\varepsilon_{2k}\|_{H^n}^2=(b-a)\lim_{k\rightarrow+\infty}\|\varepsilon_{2k}\|_{\mathbb{R}^n}^2=0.$

It is clear that \eqref{35} can be rewritten as
\begin{equation}\label{36}
\langle P_k\overline{v}_k+P_k\frac{\varepsilon_{2k}}{1+\varepsilon_{1k}}+\frac{Q_k}{1+\varepsilon_{1k}},Z_k-\overline{v}_k\rangle_{H^n}\geq 0,\quad \forall Z_k\in U^n.
\end{equation}
According to Theorem \ref{29} and Corollary \ref{37}, when $k>N$,  solving ID-VI \eqref{36} is equivalent to solving the following projection equation with the constant 
\begin{equation}\label{38}
\overline{v}_k=\mathcal{P}_{U^n}[(\mathbb{I}_{H^n}-\epsilon_0P_k)\overline{v}_k-\epsilon_0(P_k\frac{\varepsilon_{2k}}{1+\varepsilon_{1k}}+\frac{Q_k}{1+\varepsilon_{1k}})].
\end{equation}
Here  $\epsilon_0\in(0,[\epsilon_P\|P\|_{\mathcal{L}(H^n)}^{-2}]\wedge [\epsilon_P\max_{k}\|P_k\|_{\mathcal{L}(H^n)}^{-2}])\subseteq(0,2\epsilon_P\|P\|_{\mathcal{L}(H^n)}^{-2})$ can be chosen the same in \eqref{17} and \eqref{38}. We claim that the sequence $\{\|\overline{v}_k\|_{H^n}\}_{k=1}^{+\infty}$ is bounded. Indeed, we have from Lemma \ref{33} that
$\{\|P_k\|_{\mathcal{L}(H^n)}\}_{k=1}^{+\infty}$ and $\{\|Q_k\|_{H^n}\}_{k=1}^{+\infty}$ are bounded. By choosing $\delta'>0$ with $\delta_1=(1+\delta')(1-\epsilon_0\epsilon_P)<1$ and observing the inequality 
\begin{equation}\label{41}
\|c+d\|_{H^n}^2=\|c\|_{H^n}^2+\|d\|_{H^n}^2+2\langle c,d\rangle_{H^n}\leq(1+\rho)\|c\|_{H^n}^2+(1+\frac{1}{\rho})\|d\|_{H^n}^2
\end{equation}
for arbitrary $\rho>0$, we have
\begin{align*}
\|\overline{v}_k\|_{H^n}^2&=\|\mathcal{P}_{U^n}[(\mathbb{I}_{H^n}-\epsilon_0P_k)\overline{v}_k-\epsilon_0(P_k\frac{\varepsilon_{2k}}{1+\varepsilon_{1k}}+\frac{Q_k}{1+\varepsilon_{1k}})]\|_{H^n}^2\\
&\le \delta_1\|\overline{v}_k\|_{H^n}^2+\frac{2\epsilon_0^2(1+\delta')}{\delta'(1+\varepsilon_{1k})^2}(\|P_k\|_{\mathcal{L}(H^n)}^2
\|\varepsilon_{2k}\|_{H^n}^2+\|Q_k\|_{H^n}^2),
\end{align*}
where we have used the non-expansive property of metric projection operators and the fact that
$
\|(\mathbb{I}_{H^n}-\epsilon_0P_k)\overline{v}_k\|_{H^n}^2
\leq (1-\epsilon_0\epsilon_P))\|\overline{v}_k\|_{H^n}^2$ when $\epsilon_0\leq \epsilon_P\max_{k}\|P_k\|_{\mathcal{L}(H^n)}^{-2}$.
Hence
$$
\|\overline{v}_k\|_{H^n}^2\leq \max_{k}\frac{2\epsilon_0^2(1+\delta')}{\delta'(1-\delta_1)(1+\varepsilon_{1k})^2}[\max_{k}\|P_k\|_{\mathcal{L}(H^n)}^2\max_{k}\|\varepsilon_{2k}\|_{H^n}^2+\max_{k}\|Q_k\|_{H^n}^2],
$$
and the boundedness of $\{\|\overline{v}_k\|_{H^n}\}_{k=1}^{+\infty}$ follows.

Next, by choosing $\delta>0$ with $\delta_2=(1+\delta)(1-\epsilon_0\epsilon_P)<1$, from \eqref{41} one has
\begin{align*}
\|\widehat{u}-\overline{v}_k\|_{H^n}^2
\leq & \delta_2\|\widehat{u}-\overline{v}_k\|_{H^n}^2+4\epsilon^2_0(1+\frac{1}{\delta})\|P-P_k\|_{\mathcal{L}(H^n)}^2\|\overline{v}_k\|_{H^n}^2\\
&+\frac{4\epsilon^2_0(1+\frac{1}{\delta})}{(1+\varepsilon_{1k})^2}[\|P_k\|_{\mathcal{L}(H^n)}^2\|\varepsilon_{2k}\|_{H^n}^2+\|Q_k-Q\|_{H^n}^2+\varepsilon^2_{1k}\|Q\|_{H^n}^2],
\end{align*}
where we have also used the non-expansive property of metric projection operators and the fact that 
$\|(\mathbb{I}_{H^n}-\epsilon_0P_k)\overline{v}_k\|_{H^n}^2
=\|\overline{v}_k\|_{H^n}^2-2\epsilon_0\langle P_k\overline{v}_k,\overline{v}_k\rangle_{H^n}+
\epsilon_0^2\|P_k\overline{v}_k\|_{H^n}^2
\leq (1-\epsilon_0\epsilon_P)\|\overline{v}_k\|_{H^n}^2.
$
 Applying Lemma \ref{33} and the boundedness of $\{\|\overline{v}_k\|_{H^n}\}_{k=1}^{+\infty}$ gives
$\|\widehat{u}-\overline{v}_k\|_{H^n}^2\leq C_k$,
 where $\lim_{k\rightarrow+\infty}C_k=0$.

On the other hand, we have from \eqref{50} that
$
\|\widehat{u}-\overline{v}_k\|_{H^n}
\geq\|\widehat{u}-\widehat{v}_k\|_{H^n}-\frac{\varepsilon_{1k}}{1+\varepsilon_{1k}}\|\widehat{v}_k\|_{H^n}
-\|\frac{\varepsilon_{2k}}{1+\varepsilon_{1k}}\|_{H^n}.
$
Thus, $\|\widehat{u}-\widehat{v}_k\|_{H^n}\leq\frac{\varepsilon_{1k}}{1+\varepsilon_{1k}}\|\widehat{v}_k\|_{H^n}
+\frac{1}{1+\varepsilon_{1k}}\left\|\varepsilon_{2k}\right\|_{H^n}+\sqrt{C_k}.$
Letting $k\rightarrow+\infty$ on both sides gives $\lim_{k\rightarrow+\infty}\|\widehat{u}-\widehat{v}_k\|_{H^n}=0$.

For the OWAP, using Lemma \ref{33} and the boundedness of $\{\|\overline{v}_k\|_{H^n}\}_{k=1}^{+\infty}$, we have
\begin{align}\label{42}
&\lim_{k\rightarrow+\infty}\|\widehat{V}-\widehat{W}_k\|_{H^n}
= \lim_{k\rightarrow+\infty}\|P\widehat{u}+Q-P_k\widehat{v}_k-Q_k\|_{H^n}\\   
   \leq&\lim_{k\rightarrow+\infty}(\|P\|_{\mathcal{L}(H^n)}\|\widehat{u}-\widehat{v}_k\|_{H^n}+\|P-P_k\|_{\mathcal{L}(H^n)}\|\widehat{v}_k\|_{H^n}
   +\|Q-Q_k\|_{H^n})=0.\nonumber
\end{align}
It is clearly that
$$
|\langle \widehat{u},\widehat{V}\rangle_H-\langle \widehat{v}_k,\widehat{W}_k\rangle_H|
\leq\|\widehat{u}-\widehat{v}_k\|_{H^n}\|\widehat{V}\|_{H^n}+\|\widehat{v}_k\|_{H^n}\|\widehat{V}-\widehat{W}_k\|_{H^n}.
$$
Taking limit on both sides and by Theorem \ref{47}, $\lim_{k\rightarrow+\infty}|\langle \widehat{u},\widehat{V}\rangle_H-\langle \widehat{v}_k,\widehat{W}_k\rangle_H|\leq 0,$
where we have used the  boundedness of $\{\|\overline{v}_k\|_{H^n}\}_{k=1}^{+\infty}$. Thus the result follows.
\end{proof}

\end{document}
